% !TEX root = main.tex
\section{Methods}
\label{sec:methodology}\label{sec:methods}

The study has two components. A structured search identifies published retail
comparisons and supplies results from public benchmark suites. New experiments
then compare foundation models and conventional baselines under a common
protocol. The public results provide breadth, while the experiments allow
forecasts to be paired on the same observations.

\subsection{Literature search and data extraction}
\label{subsec:search}

We searched Google Scholar, Semantic Scholar, and arXiv for work published
between January 2020 and 12 April 2026. Queries combined the names of
foundation models or the terms ``foundation model'', ``pretrained model'', and
``zero-shot forecasting'' with terms for demand forecasting and benchmark
evaluation. Scopus and Web of Science were also considered, but their search
exports were not retained and are not included in the reported flow. Full
search strings and screening decisions are provided in Appendix~\ref{app:search}.

Studies were eligible if they reported quantitative forecasting results for a
foundation model on a retail dataset or on the retail portion of a public
benchmark. We excluded studies without numerical results, studies confined to
other domains, architecture papers without forecasting evaluation, and
duplicate versions of the same results. Screening and extraction were
performed by one reviewer. The search should therefore be read as a structured
map of available benchmark evidence rather than as a registered systematic
review.

For each eligible result we recorded the dataset, task, model, model family,
forecast horizon, evaluation protocol, accuracy measure, reported value,
covariate use, and available information about runtime and hardware. The
resulting file contains 802 rows, including task-level results re-extracted
from fev-bench and GIFT-Eval. Raw names were harmonised into common dataset and
model-family categories, while the original values were retained for
traceability. Table~\ref{tab:extraction-schema} summarises the fields used in
the analysis; the full data dictionary accompanies the replication files.
\label{subsec:extraction}\label{sec:extraction-schema}

\begin{table}[htbp]
\centering
\caption{Information recorded for each published benchmark result.}
\label{tab:extraction-schema}
\small
\begin{tabularx}{\textwidth}{l X}
\toprule
Group & Variables \\
\midrule
Source & Study, year, venue, benchmark suite \\
Task & Dataset, variant, frequency, number of series, forecast horizon \\
Model & Reported name, harmonised family, zero-shot or fine-tuned status \\
Evaluation & Forecast design, metric, reported value, covariate availability \\
Computation & Runtime, hardware and GPU hours when reported \\
\bottomrule
\end{tabularx}
\end{table}

\subsection{Matched retail experiments}
\label{subsec:gapfilling}

The literature search showed that standalone model papers seldom report
foundation models and gradient-boosted trees on the same retail observations.
We therefore ran both families on M5, Rohlik~v2, and Favorita. For each
dataset, the main panel contains the same 100 high-volume series for every
model. Forecasts are generated at horizons 7, 14, and 28 from common
rolling-origin windows. The experiment writes one prediction for each series,
origin, target date, and forecast step, which permits paired resampling.

The models are Seasonal Naive; recursive and direct LightGBM specifications,
including tuned variants; Chronos-Bolt-Tiny; Chronos-2;
Moirai~2.0-Small; TiRex; and TimesFM~2.5. A horizon-scaled direct LightGBM
variant is included to check whether the direct strategy benefits merely from
receiving a larger total tuning budget. M5 receives two further checks: a
sample stratified by volume and zero frequency, and a 200-trial LightGBM
tuning run. Section~\ref{sec:experimental-setup} gives the datasets, features,
and computational details.

Earlier local runs used unequal numbers of series across models. They remain
in the replication archive as provenance and for descriptive runtime
accounting, but they are not used for paired accuracy comparisons.

\subsection{Comparisons and uncertainty}
\label{subsec:metaspec}

For a foundation model $m$ and a conventional baseline $b$ evaluated in the
same cell, we report the log error ratio
\begin{equation}
  \ell=\log\left(\frac{E_m}{E_b}\right),
\end{equation}
where $E$ is the relevant error measure. Negative values favour the foundation
model. In the matched experiments, $b$ is the lowest-error baseline observed
for that dataset and horizon. Series-level bootstrap samples are shared across
models so that uncertainty reflects the paired design. Because the baseline is
selected from several candidates, the intervals describe the contrast with
the selected baseline and do not include uncertainty from that selection.

For the public benchmark analysis, a foundation-model result is paired with
the best conventional result for the same task and metric. This gives 498
comparisons: 166 each for \WAPE{}, \MASE{}, and Scaled Quantile Loss. We
summarise their medians, quantiles, and proportions below zero by benchmark
suite and metric. The published benchmark tables do not provide paired
prediction errors, so their sample sizes cannot be converted into valid
sampling variances for these log ratios. We consequently do not estimate a
pooled meta-analytic effect.

Sensitivity analyses replace the best baseline with the mean conventional
baseline, omit one suite at a time, and report metrics separately. Dataset,
horizon, and covariate availability are inspected as descriptive moderators.
With only two public suites and a limited number of independent retail
datasets, these comparisons are used to identify patterns for further testing,
not causal moderator effects. The analysis is implemented in
\texttt{analysis/meta\_regression.R}; generated tables and figures are included
with the replication materials.
